% ==============================================================================
% CONFIGURATION PREAMBLE: BÁO CÁO ĐỒ ÁN TỐT NGHIỆP VERISCHOLAR
% Tuân thủ 100% quy chuẩn Huong-dan-viet-DATN.md (HUMG)
% ==============================================================================

% 1. Thiết lập khổ giấy và căn lề trang chuẩn xác (HUMG: 2.5 - 2.5 - 3.5 - 2.5 cm)
\usepackage[
    a4paper,
    top=2.5cm,
    bottom=2.5cm,
    left=3.5cm,
    right=2.5cm,
    headheight=20pt
]{geometry}

% 2. Chuẩn hóa cỡ chữ 13pt (Word standard) và dãn dòng 1.5 lines
\usepackage[fontsize=13pt]{scrextend}
\usepackage{setspace}
\onehalfspacing

% 3. Phông chữ XeLaTeX hỗ trợ tiếng Việt Unicode bản địa
\usepackage{fontspec}
\setmainfont{Liberation Serif}[
    Path = /usr/share/fonts/truetype/liberation/,
    UprightFont = LiberationSerif-Regular.ttf,
    BoldFont = LiberationSerif-Bold.ttf,
    ItalicFont = LiberationSerif-Italic.ttf,
    BoldItalicFont = LiberationSerif-BoldItalic.ttf
]
\usepackage[provide=*,vietnamese]{babel}
\addto\captionsvietnamese{%
    \renewcommand{\listtablename}{DANH MỤC BẢNG BIỂU}%
    \renewcommand{\listfigurename}{DANH MỤC HÌNH ẢNH}%
    \renewcommand{\bibname}{TÀI LIỆU THAM KHẢO}%
}

% 4. Các gói toán học, bảng biểu và đồ họa
\usepackage{amsmath,amsfonts,amssymb}
\usepackage{graphicx}
\usepackage{booktabs}
\usepackage{tabularx}
\usepackage{longtable}
\usepackage{multirow}
\usepackage{array}
\usepackage{enumitem}

% 5. Cấu hình tiêu đề bảng biểu và hình ảnh (HUMG: Bảng ở TRÊN, Hình ở DƯỚI)
\usepackage{caption}
\captionsetup[table]{
    position=top,
    labelfont=bf,
    textfont=it,
    skip=6pt,
    name={Bảng}
}
\captionsetup[figure]{
    position=bottom,
    labelfont=bf,
    textfont=it,
    skip=6pt,
    name={Hình}
}

% 6. Cấu hình tiêu đề chân trang (Footer) và đầu trang (Header) với đường kẻ đôi chuẩn HUMG
\usepackage{fancyhdr}
\pagestyle{fancy}
\fancyhf{} % Xóa sạch cấu hình mặc định
\fancyhead[L]{\fontsize{10}{12}\selectfont\textbf{Đồ án tốt nghiệp ngành Công nghệ thông tin}}
\renewcommand{\headrule}{%
    \hrule width \headwidth height 1.2pt \vspace{1.5pt}%
    \hrule width \headwidth height 0.4pt%
}
\fancyfoot[L]{\fontsize{10}{12}\selectfont\textit{SVTH: Đàm Mạnh Dũng - 2221050359 -- Lớp: DCCTCT67\_09B}}
\fancyfoot[R]{\fontsize{10}{12}\selectfont\textit{\thepage}}
\renewcommand{\footrule}{%
    \hrule width \headwidth height 0.4pt \vspace{1.5pt}%
    \hrule width \headwidth height 1.2pt%
}

% Áp dụng phong cách chuẩn cho trang mở đầu chương / plain
\fancypagestyle{plain}{
    \fancyhf{}
    \fancyhead[L]{\fontsize{10}{12}\selectfont\textbf{Đồ án tốt nghiệp ngành Công nghệ thông tin}}
    \renewcommand{\headrule}{%
        \hrule width \headwidth height 1.2pt \vspace{1.5pt}%
        \hrule width \headwidth height 0.4pt%
    }
    \fancyfoot[L]{\fontsize{10}{12}\selectfont\textit{SVTH: Đàm Mạnh Dũng - 2221050359 -- Lớp: DCCTCT67\_09B}}
    \fancyfoot[R]{\fontsize{10}{12}\selectfont\textit{\thepage}}
    \renewcommand{\footrule}{%
        \hrule width \headwidth height 0.4pt \vspace{1.5pt}%
        \hrule width \headwidth height 1.2pt%
    }
}

% 7. Định dạng Mục lục chuẩn xác 100% theo mẫu đồ án xuất sắc HUMG (Hình 2 & 3)
\usepackage[titles]{tocloft}

% Dấu chấm dẫn trang (leader dots) cho tất cả các cấp (kể cả chapter)
\renewcommand{\cftchapleader}{\cftdotfill{\cftdotsep}}
\renewcommand{\cftsecleader}{\cftdotfill{\cftdotsep}}
\renewcommand{\cftsubsecleader}{\cftdotfill{\cftdotsep}}
\renewcommand{\cftsubsubsecleader}{\cftdotfill{\cftdotsep}}

% Khoảng cách dấu chấm dày dặn (dotsep = 1) giống Word
\renewcommand{\cftdotsep}{1}

% Tiền tố "CHƯƠNG " cho các chương có đánh số
\renewcommand{\cftchappresnum}{CHƯƠNG }
\renewcommand{\cftchapaftersnum}{:\ }
\newlength{\cftchaplen}
\settowidth{\cftchaplen}{\bfseries CHƯƠNG 1:\ }
\addtolength{\cftchapnumwidth}{\cftchaplen}

% Dấu chấm sau số thứ tự section và subsection (1.1., 1.2., 2.1.1.)
\renewcommand{\cftsecaftersnum}{.}
\renewcommand{\cftsubsecaftersnum}{.}
\renewcommand{\cftsubsubsecaftersnum}{.}

% Thụt đầu dòng chuẩn xác theo phân cấp
\cftsetindents{chapter}{0pt}{3.0cm}
\cftsetindents{section}{0.5cm}{0.9cm}
\cftsetindents{subsection}{1.0cm}{1.2cm}
\cftsetindents{subsubsection}{1.5cm}{1.4cm}

% Khoảng cách đều đặn, không nhảy dòng quá xa giữa các chương
\setlength{\cftbeforechapskip}{4pt}
\setlength{\cftbeforesecskip}{2pt}
\setlength{\cftbeforesubsecskip}{1pt}

% Phông chữ trong mục lục: chuẩn mực, thanh thoát
\renewcommand{\cftchapfont}{\normalfont}
\renewcommand{\cftchappagefont}{\normalfont}
\renewcommand{\cftsecfont}{\normalfont}
\renewcommand{\cftsecpagefont}{\normalfont}

% 7. Phân cấp đề mục và độ sâu mục lục (Tối đa 4 cấp)
\usepackage{titlesec}
\setcounter{secnumdepth}{4}
\setcounter{tocdepth}{3}

% Định dạng tiêu đề chương không bị cách khoảng quá xa
\titleformat{\chapter}[display]
  {\normalfont\bfseries\centering}
  {\MakeUppercase{\chaptertitlename}\ \thechapter}{10pt}{\Large\MakeUppercase}
\titlespacing*{\chapter}{0pt}{-10pt}{20pt}

% 8. Liên kết siêu văn bản PDF học thuật
\usepackage{hyperref}
\hypersetup{
    colorlinks=true,
    linkcolor=black,
    citecolor=blue,
    urlcolor=blue,
    pdfauthor={Đàm Mạnh Dũng},
    pdftitle={Xây dựng nền tảng AI hỗ trợ nghiên cứu và viết bài báo khoa học VinUni dựa trên Agentic RAG, truy xuất tri thức học thuật và kiểm chứng trích dẫn}
}
